\clearpage
\chapter*{Abstract}
\addcontentsline{toc}{chapter}{Abstract}
\markboth{Abstract}{Abstract}

\noindent
This thesis presents the design and implementation of a real-time supervision and cascade
control system for a thermal process based on a Siemens S7-1500 programmable logic
controller. The installation under study consists of a plate heat exchanger in which a hot
fluid, pre-heated by an electric furnace, exchanges energy in counter-current with a cold
fluid circuit.

A Flask-based web application communicates with the PLC via two complementary protocols:
Snap7 for direct data block access (DB1), and the Siemens Open Pipe mechanism for reading
and writing WinCC Unified tags over a Unix domain socket. An open-loop step-response
identification experiment is conducted on the inner process
($F_{1,SP} \rightarrow T_{in1,HE1}$) to extract the parameters of a First-Order Plus Time
Delay model: static gain $K$, time constant $\tau$, and dead time $\theta$. These parameters
are used to tune a PI controller with proportional action on the measurement (PI+P) following
Internal Model Control rules.

This inner PI+P controller regulates the hot inlet temperature $T_{in1}$, forming the fast
inner loop of a cascade architecture. An outer proportional controller closes the cascade by
adjusting the $T_{in1}$ setpoint to achieve the desired cold outlet temperature $T_{out2}$.
The supervision interface runs inside a Docker container and is accessible from a PC on the
laboratory network and, via Siemens Industrial Edge, from the embedded browser of a SIMATIC
Unified 7'' touch panel, displaying a real-time P\&ID synoptic
and providing a cascade identification and commissioning tool.

\medskip
\noindent\textbf{Keywords:} PLC, S7-1500, cascade control, process identification, Flask,
Docker, WinCC Unified, Industrial Edge, heat exchanger, PI+P controller, IMC tuning

\bigskip
\hrule
\bigskip

\begin{otherlanguage}{french}
\chapter*{Résumé}
\addcontentsline{toc}{chapter}{Résumé}
\markboth{Résumé}{Résumé}

\noindent
Ce mémoire présente la conception et l'implémentation d'un système de supervision temps réel
et de contrôle en cascade d'un procédé thermique basé sur un automate Siemens S7-1500.
L'installation étudiée est un échangeur de chaleur à plaques dans lequel un fluide chaud,
préchauffé par un four électrique, échange de l'énergie à contre-courant avec un circuit
froid.

Une application web développée en Flask communique avec l'automate via deux protocoles
complémentaires~: Snap7 pour l'accès direct au bloc de données DB1, et le mécanisme Open
Pipe de Siemens pour la lecture et l'écriture des variables WinCC Unified via une socket Unix.
Une identification expérimentale par test échelon en boucle ouverte permet d'extraire les
paramètres du modèle du premier ordre avec retard (gain $K$, constante de temps $\tau$,
retard $\theta$), utilisés pour régler un régulateur PI à action proportionnelle sur la mesure
selon la méthode IMC.

Ce régulateur PI+P constitue la boucle interne rapide d'une structure en cascade, régulant
la température d'entrée chaude de l'échangeur. Une boucle externe proportionnelle pilote
la consigne de cette boucle interne pour atteindre la température de sortie froide souhaitée.
L'interface de supervision est déployée via Docker et accessible depuis un poste de
développement ou directement depuis le panel SIMATIC Unified 7'' via Industrial Edge.

\medskip
\noindent\textbf{Mots-clés~:} automate, S7-1500, contrôle en cascade, identification,
Flask, Docker, WinCC Unified, Industrial Edge, échangeur de chaleur, régulateur PI+P, IMC
\end{otherlanguage}
